\documentclass[../../../uem_phy_std_q.tex]{subfiles}

\begin{document}
    \setcounter{figure}{0}
    電流と電圧の関係が図\ref{fig:fig_uem_phy_std_em_cx_08_01_q}で表されるような電球がある．\\
    この電球と，電池（起電力$E=9.0\:\mathrm{V}$），電気抵抗（抵抗値\\
    $r=4.0\:\Omega$）を用いて回路を作った．

    \begin{enumerate}
        \item%
            図\ref{fig:fig_uem_phy_std_em_cx_08_02_q}の回路を作ったとき，電球に流れる電流$I_1$を求めよ．
        \item%
            図\ref{fig:fig_uem_phy_std_em_cx_08_03_q}の回路を作ったとき，電球に流れる電流$I_2$を求めよ．
        \item%
            図\ref{fig:fig_uem_phy_std_em_cx_08_04_q}のような2通りの回路AとBを作ったところ，%
            回路Aと回路Bの電球の明るさが同じになった．抵抗値$R_3$を求めよ．
    \end{enumerate}
    \begin{minipage}{\linewidth}
        \setlength{\abovecaptionskip}{-1pt}
        \centering
            \includegraphics{\subfix{fig_uem_phy_std_em_cx_08/fig_uem_phy_std_em_cx_08_01_q.pdf}}
            \captionof{figure}{} \label{fig:fig_uem_phy_std_em_cx_08_01_q}
            \vspace{7pt}
            \begin{minipage}{.39\linewidth}
                \setlength{\abovecaptionskip}{3pt}
                \centering
                    \includegraphics{\subfix{fig_uem_phy_std_em_cx_08/fig_uem_phy_std_em_cx_08_02_q.pdf}}
                    \captionof{figure}{} \label{fig:fig_uem_phy_std_em_cx_08_02_q}
            \end{minipage}
            \begin{minipage}{.58\linewidth}
                \setlength{\abovecaptionskip}{3pt}
                \centering
                    \includegraphics{\subfix{fig_uem_phy_std_em_cx_08/fig_uem_phy_std_em_cx_08_03_q.pdf}}
                    \captionof{figure}{} \label{fig:fig_uem_phy_std_em_cx_08_03_q}
            \end{minipage}
    \end{minipage}
    \\
    \begin{minipage}{\linewidth}
        \centering
            \includegraphics{\subfix{fig_uem_phy_std_em_cx_08/fig_uem_phy_std_em_cx_08_04_q.pdf}}
            \captionof{figure}{} \label{fig:fig_uem_phy_std_em_cx_08_04_q}
    \end{minipage}    
\end{document}



